\section{切换过程中的连续状态与初始条件}

直流稳态、开关切换瞬间与随后暂态是三种不同状态。稳态直流下，理想电感的端电压为零、可按导线处理；电容的传导电流为零、可按断路处理。这些替代只用于求稳态，不能直接搬到切换瞬间。切换时没有电压冲激的电感电流连续，没有电流冲激的电容电压连续，它们把切换前的状态传递到新电路。

分析开关题时，先画切换前的稳态等效电路，求每个电感电流与电容电压；再画切换后的真实闭合通路，保留这些初值，求新的端电压和支路电流。灯泡可能与电感同支路，也可能成为电感释放磁能的回路，不能只凭“有自感所以灯会闪”判断亮度。灯的瞬时功率取决于自身 $i^2R$，方向反转不会使亮度成为负值。

电感开路时若找不到放电通路，理想电路通常不自洽，真实装置会经寄生电容或击穿释放能量。本讲切换题均明确保留闭合通路。电容从一电源改接另一电源时也可能有快速再分配过程，若理想零电阻造成矛盾，应保留小电阻后求极限，而不是同时强加互不相容的初值。

\fig{em5:initial}{元件正向与状态连续条件；连续性来自相应电压或电流的短时积分}{\begin{tikzpicture}
\draw[world](0,2.4)--(0,1.95);\draw[world](0,0.35)arc(-90:90:.2);\draw[world](0,0.75)arc(-90:90:.2);\draw[world](0,1.15)arc(-90:90:.2);\draw[world](0,1.5500000000000003)arc(-90:90:.2);\draw[world](0,.35)--(0,0);
\draw[flow,red](0,3.1)--(0,2.5)node[midway,right]{$i_L$};\node[left]at(-.3,2.35){$+$};\node[left]at(-.3,0){$-$};\node[left]at(-.5,1.1){$u_L$};\node[right]at(.4,1.1){$L$};
\node at(0,-.7){$L\Delta i_L=\int u_L\dd t$};
\begin{scope}[xshift=5cm]\draw[world](0,0)--(0,.9) (0,1.3)--(0,2.4);\draw[world](-.5,.9)--(.5,.9) (-.5,1.3)--(.5,1.3);
\draw[flow,red](0,3.1)--(0,2.5)node[midway,right]{$i_C$};\node[left]at(-.3,2.35){$+$};\node[left]at(-.3,0){$-$};\node[left]at(-.7,1.1){$u_C$};\node[right]at(.7,1.1){$C$};\node at(0,-.7){$C\Delta u_C=\int i_C\dd t$};\end{scope}

\end{tikzpicture}}
图\ref{em5:initial}的两行分别由$L\Delta i=\int u_L\dd t$和$C\Delta u=\int i_C\dd t$得到。切换间隔趋零且相应被积量有限时，积分趋零；两种连续性应各自检查，不能由其中一种推出另一种。

\subsection{快慢阶段与冲激改变状态的情形}

若理想开关把不同电压的电容直接并接，电流中出现冲激，原来每只电容的电压就不再分别连续；能保留的是并接导体总电荷。类似地，互感中的理想电流阶跃可以给另一个电感施加电压冲激。连续性必须检查适用条件，不能只凭元件名称判断。

\fig{em5:redistribution}{电压不同的电容通过小电阻接通，电荷守恒而各自电压在快过程中改变}{\begin{tikzpicture}

\foreach \x/\j in {0/1,4/2}{\draw[world](\x,0)--(\x,.9) (\x,1.3)--(\x,2.4);\draw[world]({\x-.45},.9)--({\x+.45},.9) ({\x-.45},1.3)--({\x+.45},1.3);\node[left]at({\x-.5},1.1){$C_{\j}$};}
\draw[world](0,0)--(4,0) (0,2.4)--(1.2,2.4) (2,2.4)--(2.4,2.4) (2.9,2.4)--(4,2.4);\draw[world](1.2,2.2)rectangle(2,2.6);\node[above]at(1.6,2.65){$r\to0^+$};\fill(2.4,2.4)circle(1.5pt);\fill(2.9,2.4)circle(1.5pt);\draw[world](2.4,2.4)--(2.85,2.7);\node[above]at(2.7,2.85){$S$};
\node at(0,3.2){$U_1$};\node at(4,3.2){$U_2$};\node at(2,-.8){$U_f=(C_1U_1+C_2U_2)/(C_1+C_2)$};

\end{tikzpicture}}
图\ref{em5:redistribution}以接上侧公共导体的极板电荷为正。小电阻上的初始压差为 $U_1-U_2$，电流随 $r$ 减小而变大，重新分配时间约为 $rC_1C_2/(C_1+C_2)$。电流积分仍可有限，故不能强加每只电容电压在这个极限快过程两侧都不变。

由上侧导体的总电荷守恒先求图中的 $U_f$，再相减电场能，得
\[Q_J=\frac12C_1U_1^2+\frac12C_2U_2^2-\frac12(C_1+C_2)U_f^2
=\frac{C_1C_2}{2(C_1+C_2)}(U_1-U_2)^2.\]
即使 $r\to0^+$，热量也不趋零。这里用小电阻给出可实现的快过程极限，严格的无电阻、无电感瞬时并接模型本身不能描述能量去向。


对含有快慢两种时间尺度的电路，可以把短促的重新分配与之后的平滑演化分开。快过程中忽略那些在短时间内来不及显著改变的量，对相应微分方程积分，建立前后状态的连接关系。慢过程中使用新电路的方程，并把快过程末态作为初态。这样的处理仍需追踪快过程的热量，不能因为持续时间短就一概视为能量不变。

例如电容重新分配电荷时，细小接触电阻中的电流虽持续很短，但其平方积分可以保持有限。先用电荷守恒求末电压，再比较前后电能，通常比试图求无限窄电流峰更可靠。电感切换到带限压元件的续流路径时，则应先按元件导通条件求电流降至零的事件时刻，此后回路方程已经改变，不能把衰减解延伸成反向电流。


\section{一阶电路与时间常量}

理想电感 $L$ 串联电阻 $R$ 接到恒压 $E$ 上，回路有 $L\dot i+Ri=E$。令最终电流 $I_\infty=E/R$，减去稳态关系后得 $L\dd(i-I_\infty)/\dd t=-R(i-I_\infty)$。令 $\delta=i-I_\infty$。若 $\delta(0)\ne0$，分离变量并保留积分上下限，
\[
\int_{\delta(0)}^{\delta(t)}\frac{\dd z}{z}
=-\frac RL\int_0^t\dd s,\qquad
\ln\left|\frac{\delta(t)}{\delta(0)}\right|=-\frac{Rt}{L}.
\]
偏差的符号不会在有限时间内改变，故 $\delta(t)=\delta(0)\e^{-Rt/L}$；若初始偏差为零，原方程直接给出恒定的稳态解。代入 $\delta(0)=I_0-I_\infty$，得到
\[i=I_\infty+(I_0-I_\infty)\e^{-t/\tau_L},\qquad\tau_L=L/R.\]
电感电压为 $u_L=E-Ri$，自感电动势为 $-u_L$。电源断开并经电阻续流时，$E=0$，电流按相同形式衰减。$t=\tau_L$ 时剩余偏差为原来的 $\e^{-1}$，并不是过程已经结束；指定误差 $\eta$ 时应取 $t\ge\tau_L\ln(1/\eta)$。

电容充电回路满足 $R\dot q+q/C=E$。令 $z=q-CE$，则 $\dot z=-z/(RC)$，于是
\[
\int_{q_0-CE}^{q-CE}\frac{\dd z}{z}
=-\int_0^t\frac{\dd s}{RC},\qquad
q-CE=(q_0-CE)\e^{-t/(RC)}.
\]
初始偏差为零时仍应单独取恒定解；其他初态给出 $q=CE+(q_0-CE)\e^{-t/(RC)}$。时间常量 $RC$ 与电感的 $L/R$ 对电阻的依赖相反：电阻越大，电容充电越慢，电感电流建立相对最终值则越快。这里比较的是归一化时间过程，电感的最终电流 $E/R$ 本身还会随电阻改变。

\fig{em5:time-constant}{一阶过程的剩余偏差；虚线为初始切线及一个时间常量的读数}{
\begin{tikzpicture}[x=1.5cm,y=2.2cm]
\draw[axis](0,0)--(4.4,0)node[right]{$t/\tau$};
\draw[axis](0,0)--(0,1.25)node[above]{$\delta/\delta_0$};
\draw[world,blue,domain=0:4.1,samples=80]plot(\x,{exp(-\x)});
\draw[guide](0,1)--(1,0);
\draw[guide](1,0)--(1,{exp(-1)})--(0,{exp(-1)});
\node[left]at(0,1){$1$};
\node[left]at(0,{exp(-1)}){$\e^{-1}$};
\node[below]at(1,0){$1$};
\node at(2.6,.6){$\delta=\delta_0\e^{-t/\tau}$};
\end{tikzpicture}}

图\ref{em5:time-constant}把电感电流偏差$\delta=i-I_\infty$或电容电压偏差$\delta=u-U_\infty$除以各自初始偏差$\delta_0$。曲线在$t=\tau$处的高度为$\e^{-1}$，初始切线则在同一横坐标截到零；实际曲线始终不在有限时间到零。初值高于或低于新稳态时，$\delta_0$符号不同，但归一化后的曲线相同。$\delta_0=0$时系统从开始就在稳态，无须作此归一化。


复杂的线性电阻网络接单个电感时，可先对电感端口求戴维南等效电压 $E_{\rm th}$ 和电阻 $R_{\rm th}$，再使用 $L/R_{\rm th}$。求等效电阻时独立理想电压源短接、理想电流源断开，电源内阻仍保留。开关改变后要重新求等效，不能使用切换前的电阻值。


\subsection{戴维南等效与非零初态}

对只有一个电感的线性网络，把其他元件等效到其两端以后，可将电感电流写成“新稳态加衰减偏差”。这个形式不要求初电流为零，初值高于新稳态时电流下降，自感电动势顺原电流；初值低于新稳态时电流上升，自感作用反向。两种情况的电流都连续，但电感端电压可以在开关切换后立刻改变。

若电路中同时有电感和电容，通常需要两个初始状态。仅知道切换前端电压，不足以确定切换后全部运动；还要知道电感中存储的电流。把未知初值误当成零，会改变振幅和相位，甚至改变电容最先带电的极性。对电容并联电感的释放过程，最可靠的方法是先写节点电流守恒，再根据它确定电容电压的初始导数。

实际开关存在断开再接通的微小时间。若这段时间中电感没有续流路径，端电压可能由寄生电容决定，之后再接通的初态与理想无缝换接不同。本讲采用理想快速换接且续流通路明确的条件；解答中的连续性是相应有限电压积分的结果，不是脱离连接关系的口号。

\fig{em5:thevenin-network}{单电感端口的电阻网络及其等效回路；初始电流由切换前的电路另行确定}{\begin{tikzpicture}

\draw[world](0,0)--(4,0) (0,2.4)--(.9,2.4) (1.7,2.4)--(4,2.4) (0,0)--(0,.85) (0,1.55)--(0,2.4);
\draw[world](0,1.2)circle(.35);\node at(0,1.35){$+$};\node at(0,1){$-$};\node[left]at(-.4,1.2){$E$};\draw[world](.9,2.2)rectangle(1.7,2.6);\node[above]at(1.3,2.65){$R_s$};
\draw[world](2.6,0)--(2.6,.8) (2.6,1.5)--(2.6,2.4);\draw[world](2.4,.8)rectangle(2.8,1.5);\node[left]at(2.3,1.15){$R_p$};
\draw[world](4,2.4)--(4,1.95);\draw[world](4,0.35)arc(-90:90:.2);\draw[world](4,0.75)arc(-90:90:.2);\draw[world](4,1.15)arc(-90:90:.2);\draw[world](4,1.5500000000000003)arc(-90:90:.2);\draw[world](4,.35)--(4,0);
\node[right]at(4.4,1.1){$L$};\node[right]at(4.1,2.4){$A$};\node[right]at(4.1,0){$B$};
\begin{scope}[xshift=7cm]
\draw[world](0,0)--(3,0) (0,2.4)--(.8,2.4) (1.6,2.4)--(3,2.4) (0,0)--(0,.85) (0,1.55)--(0,2.4);\draw[world](0,1.2)circle(.35);\node at(0,1.35){$+$};\node at(0,1){$-$};\node[left]at(-.4,1.2){$E_{\rm th}$};\draw[world](.8,2.2)rectangle(1.6,2.6);\node[above]at(1.2,2.65){$R_{\rm th}$};
\draw[world](3,2.4)--(3,1.95);\draw[world](3,0.35)arc(-90:90:.2);\draw[world](3,0.75)arc(-90:90:.2);\draw[world](3,1.15)arc(-90:90:.2);\draw[world](3,1.5500000000000003)arc(-90:90:.2);\draw[world](3,.35)--(3,0);
\node[right]at(3.4,1.1){$L$};\draw[flow,red](2,2.4)--(2.7,2.4)node[midway,above]{$i$};\end{scope}

\end{tikzpicture}}
以图\ref{em5:thevenin-network}为例，将电感移开后由分压得 $E_{\rm th}=ER_p/(R_s+R_p)$；再将独立理想电压源置零，得到 $R_{\rm th}=R_s\parallel R_p$。于是
\[i(t)=\frac{E_{\rm th}}{R_{\rm th}}+
\left(I_0-\frac{E_{\rm th}}{R_{\rm th}}\right)\e^{-R_{\rm th}t/L}.\]
初流 $I_0$ 是换接前的电感状态，既不能擅自置零，也不能用新稳态替代。右图只等效电感端口，不保证原网络每只电阻的电流、功率均等于某个等效元件中的量；需要这些内部量时应代回左图。


对电阻网络求端口等效，还可用短路电流检验：$I_{\rm sc}=E_{\rm th}/R_{\rm th}$ 应等于把电感替换为理想导线时的稳态电流。双电源供电题中，初流是两路支流之和，换接后应按新的电源和负载求新稳态。电流初值与新稳态不是同一次等效计算的结果。


\subsection{阶跃换接、有限脉冲与误差要求}

\fig{em5:rc-switch}{电容从一个恒压源无缝换接到另一个恒压源；电压按上极板对下极板计}{
\begin{tikzpicture}[scale=.86]
\foreach \shift/\src/\state in {0/1/切换前稳态,5.8/2/切换后回路}{
\begin{scope}[xshift=\shift cm]
\node at(1.65,3){\state};
\draw[world](0,0)--(3.3,0)--(3.3,.9)
 (3.3,1.3)--(3.3,2.3)--(2.2,2.3)
 (1.2,2.3)--(0,2.3)--(0,1.5) (0,.7)--(0,0);
\draw[world](0,1.1)circle(.4);
\node at(0,1.1){$E_{\src}$};
\node[left]at(-.45,1.5){$+$};
\draw[world](1.2,2.12)rectangle(2.2,2.48);
\node[above]at(1.7,2.5){$R$};
\draw[world](3,.9)--(3.6,.9) (3,1.3)--(3.6,1.3);
\node[right]at(3.65,1.1){$C$};
\node[left]at(2.95,1.5){$+$};
\draw[flow](2.4,2.3)--(3,2.3)node[midway,above]{$i$};
\end{scope}}
\node[align=center,text width=4.6cm]at(1.65,-.85){$u_C(0^-)=E_1$\\$i(0^-)=0$};
\node[align=center,text width=5.2cm]at(7.45,-.85){$u_C(0^+)=E_1$\\$i(0^+)=(E_2-E_1)/R$};
\end{tikzpicture}}

图\ref{em5:rc-switch}左图先用稳态电流为零确定电容原有电压，右图保留这一电压、换入新电源。沿右图正电流方向列式，得到$E_2-Ri-u_C=0$，再用$i=C\dot u_C$建立下面的暂态方程。电容电压连续，电阻电流却可以立即改变；当$E_2<E_1$时，图中电流的计算值为负，表示实际电流反向。新稳态是$u_C(\infty)=E_2$，不能用它代替初值。


电容在电源从 $E_1$ 改为 $E_2$ 的阶跃中，若切换前已稳态，初电压 $U_1=E_1$，之后满足 $RC\dot u+u=E_2$。令 $w=u-E_2$，则 $w(0)=E_1-E_2$。当 $E_1\ne E_2$ 时，分离变量并按初值确定积分下限：
\begin{align*}
\frac{\dd w}{w}&=-\frac{\dd t}{RC},\\
\int_{E_1-E_2}^{w(t)}\frac{\dd z}{z}
&=-\frac1{RC}\int_0^t\dd s,\\
\left[\ln|z|\right]_{E_1-E_2}^{w(t)}
&=-\frac{t}{RC},\\
\ln\left|\frac{w(t)}{E_1-E_2}\right|
&=-\frac{t}{RC}.
\end{align*}
取指数并由初值确定符号，得到 $w(t)=(E_1-E_2)\e^{-t/(RC)}$。
若 $E_1=E_2$，原方程的初始状态已经是平衡态，故 $w(t)=0$，同一结果仍成立。
再对电压求导，
$C\dot u=C(E_1-E_2)(-1/RC)\e^{-t/(RC)}$，故
\[u(t)=E_2+(E_1-E_2)\e^{-t/(RC)},\quad i(t)=\frac{E_2-E_1}{R}\e^{-t/(RC)}.\]
当 $E_2<E_1$ 时，电流反向，电容向网络释放能量；外电源也可能吸收能量。热量由电流平方积分求出：
\begin{align*}
Q_J&=\frac{(E_2-E_1)^2}{R}\int_0^\infty\e^{-2t/(RC)}\,\dd t\\
&=\frac{(E_2-E_1)^2}{R}
\left[-\frac{RC}{2}\e^{-2t/(RC)}\right]_0^\infty
=\frac C2(E_2-E_1)^2.
\end{align*}
指数求导带出 $-2/(RC)$，所以原函数须乘 $-RC/2$。不能只用电容初末能量差，因为电源在全过程还做有向功 $E_2C(E_2-E_1)$。

串联 RL 的电源也可在任意时刻改变。每个恒压阶段的解由该阶段稳态与初始偏差组成，在切换点保持电感电流相等即可接续。对于短脉冲，脉冲结束时电流一般尚未到稳态，续流释放的磁能要用这个实际末电流计算，而不是用电源电压除以电阻的稳态值。

\fig{em5:pulse-current}{有限恒压脉冲后经同一电阻续流；换接电流连续，续流初值不是稳态值}{\begin{tikzpicture}

\draw[axis](-.2,0)--(6.2,0)node[right]{$t/\tau$};\draw[axis](0,-.2)--(0,2.6)node[above]{$i/(E/R)$};
\draw[world,blue,domain=0:2,samples=60]plot(\x,{2*(1-exp(-\x))});\draw[world,red,domain=2:5.6,samples=70]plot(\x,{2*(1-exp(-2))*exp(-(\x-2))});
\draw[guide](0,2)--(5.6,2);\draw[guide](2,0)--(2,{2*(1-exp(-2))});\node[left]at(0,2){$1$};\node[below]at(2,0){$T/\tau=2$};\node at(4,1.5){$i(T)\e^{-(t-T)/\tau}$};

\end{tikzpicture}}
图\ref{em5:pulse-current}示 $T=2\tau$ 的脉冲，换接时 $i(T)=(E/R)(1-\e^{-T/\tau})$。后半段只改变方程的外源项，不能改变这个状态；续流全过程发热为 $Li(T)^2/2$。若续流电阻另有数值，后半段还须使用新的时间常量。


“五个时间常量后稳定”是一种误差约定。指数偏差在 $5\tau$ 时仍有约0.67\%，若要求千分之一误差，应使用 $\ln1000\simeq6.91$ 个时间常量。题目没有规定精度时，严格数学解用极限表示稳态，工程估计则需要给出所采用的容差。

时间常量还应与测量时段比较。若观测时间远小于 $L/R$，电感电流初期近似 $i\simeq Et/L$；若远小于 $RC$，未充电电容电压近似 $u\simeq Et/(RC)$。这些展开来自 $\e^{-z}\simeq1-z$，只有 $|z|\ll1$ 才成立，不能延伸到整个充放电过程。


\section{总电量、有限暂态量与能量}

电流趋于非零常量时，从接通一直积分到无穷时刻的总电量必然发散。电路“稳定”不是严格的有限终点；若题目问总电量，应给有限观察时间，或明确只问暂态偏差对应的电量。对于从零建立的电感电流，
\[
Q(T)=\frac ER\left[t+\tau_L\e^{-t/\tau_L}\right]_0^T
=\frac ER[T-\tau_L(1-\e^{-T/\tau_L})].
\]
这里 $\dd(\tau_L\e^{-t/\tau_L})/\dd t=-\e^{-t/\tau_L}$，原函数在下限 $0$ 的值为 $\tau_L$，不能漏减。
相对于从一开始就达到稳态而少通过的电量，则为
\[
\int_0^\infty(I_\infty-i)\dd t
=I_\infty\left[-\tau_L\e^{-t/\tau_L}\right]_0^\infty
=I_\infty\tau_L=\frac{EL}{R^2}.
\]这两种问题的积分对象不同，不能把有限的暂态欠额当作无限时间总电量。

\fig{em5:charge-areas}{同一建立电流曲线下的有限时间电量与上方的暂态欠额，是不同的积分区域}{\begin{tikzpicture}

\draw[axis](-.2,0)--(5.6,0)node[right]{$t/\tau_L$};\draw[axis](0,-.2)--(0,2.6)node[above]{$i/I_\infty$};
\fill[blue!15](0,0)--plot[domain=0:2,samples=50](\x,{2*(1-exp(-\x))})--(2,0)--cycle;
\fill[red!12](0,2)--(5,2)--plot[domain=5:0,samples=90](\x,{2*(1-exp(-\x))})--cycle;
\draw[world,blue,domain=0:5,samples=80]plot(\x,{2*(1-exp(-\x))});\draw[guide](0,2)--(5.1,2);\draw[guide](2,0)--(2,{2*(1-exp(-2))});
\node[left]at(0,2){$1$};\node[below]at(2,0){$T/\tau_L$};\node at(1.3,.55){$Q(T)$};\node at(3,3.3){$\Delta Q$：上方面积延伸至无穷};

\end{tikzpicture}}
图\ref{em5:charge-areas}中的蓝色区域只积到 $T$，红色区域应沿水平渐近线一直延伸。若把两个区域都称为“总电量”，便会把一个随观察时间增长的量与一个有限的暂态差额混淆。


在电容充电过程中，电源做功 $E\Delta q$，电容能增量由 $q^2/(2C)$ 的始末差确定，差额为热。未充电电容接恒压源的全过程中，电源做功 $CE^2$，电容只储存 $CE^2/2$，另有一半耗散；这一结论与充电电阻大小无关，但充电时间有关。电感续流释放的总热则为初始磁场能 $LI_0^2/2$，前提是没有其他储能或外源继续做功。


\fig{em5:rc-energy}{未充电电容接恒压源时，累计电源功、储能和热量的关系}{\begin{tikzpicture}

\draw[axis](-.2,0)--(5.5,0)node[right]{$t/(RC)$};\draw[axis](0,-.2)--(0,2.7)node[above]{$W/(CE^2)$};
\draw[black,thick,domain=0:5,samples=80]plot(\x,{2*(1-exp(-\x))});\draw[blue,thick,domain=0:5,samples=80]plot(\x,{(1-exp(-\x))^2});\draw[red,thick,domain=0:5,samples=80]plot(\x,{1-exp(-2*\x)});
\node[right]at(5,2){电源功};\node[red]at(2.5,1.35){焦耳热};\node[blue]at(3,.6){电场能};\node[left]at(0,2){$1$};\node[left]at(0,1){$1/2$};\draw[guide](0,1)--(5,1);

\end{tikzpicture}}
令 $s=t/(RC)$。图\ref{em5:rc-energy}三条曲线分别为 $W_s/(CE^2)=1-\e^{-s}$、$W_E/(CE^2)=(1-\e^{-s})^2/2$、$Q_J/(CE^2)=(1-\e^{-2s})/2$，任意时刻均有 $W_s=W_E+Q_J$。改变电阻只改横轴时间尺度；保留初态和末态后，全过程能量比例不变。


\section{LC振荡的方程、初相位与能量交换}

在无电阻的闭合 LC 回路，约定电容上极板电荷为 $q$，电流 $i=-\dot q$ 表示上极板放电。回路电压关系为 $L\dot i=q/C$，消去电流得 $\ddot q+q/(LC)=0$。令 $\omega_0=1/\sqrt{LC}$，试取 $q=\e^{st}$，则 $\dot q=s\e^{st}$、$\ddot q=s^2\e^{st}$。代入并约去指数得 $s^2+\omega_0^2=0$，两个根为 $s=\pm\mathrm j\omega_0$；将共轭指数解合成实函数，一般解为
\[q=A\cos\omega_0t+B\sin\omega_0t.\]
两个常数分别由初始电荷和初始电流决定。若 $q(0)=Q_0,i(0)=I_0$，则 $A=Q_0,B=-I_0/\omega_0$，因而 $i=Q_0\omega_0\sin\omega_0t+I_0\cos\omega_0t$。周期 $T_0=2\pi\sqrt{LC}$，电流与电荷相差四分之一周期，但谁先达到峰值须结合符号与初态判断。

两端能量为 $W_E=q^2/(2C)$、$W_B=Li^2/2$。逐项使用复合函数求导，
\[
\dot W_E=\frac qC\dot q,\qquad
\dot W_B=Li\dot i=-L\dot q\,\dot i.
\]
所以 $\dot W_E+\dot W_B=\dot q(q/C-L\dot i)=0$，总能量守恒。电容电压为零时电流绝对值最大；电流为零时电荷绝对值最大。四分之一周期后仅仅是能量完成一次单向转移，半周期后电容极性反转，一整周期后所有有向状态量才回到初值。

\fig{em5:lc}{电流从带电极板流出后进入电感；$i=-\dot q$，电感被动端压为 $q/C$}{\begin{tikzpicture}
\draw[world](3.5,2.4)--(3.5,1.95);\draw[world](3.5,0.35)arc(-90:90:.2);\draw[world](3.5,0.75)arc(-90:90:.2);\draw[world](3.5,1.15)arc(-90:90:.2);\draw[world](3.5,1.5500000000000003)arc(-90:90:.2);\draw[world](3.5,.35)--(3.5,0);
\draw[world](0,0)--(3.5,0) (0,2.4)--(3.5,2.4) (0,0)--(0,.9) (0,1.3)--(0,2.4);\draw[world](-.45,.9)--(.45,.9) (-.45,1.3)--(.45,1.3);
\node[left]at(-.55,1.1){$C$};\node[right]at(.5,1.5){$+q$};\node[right]at(.5,.6){$-q$};\node[right]at(3.9,1.1){$L$};\node[right]at(3.6,2.25){$+$};\node[right]at(3.6,.1){$-$};\draw[flow,red](1,2.4)--(2.4,2.4)node[midway,above]{$i$};

\end{tikzpicture}}
图\ref{em5:lc}的电流箭头从电容上极板流出，所以该极板电荷减少时$i>0$，即$i=-\dot q$。沿图示上支路向右经过电感，电容电压推动电流建立，给出$L\dot i=q/C$。后续相位图和时间图沿用这一正向。


\section{串联RLC的阻尼与能量耗散}

\fig{em5:rlc-circuit}{串联RLC自由回路沿放电电流列式，电阻持续耗散两种储能}{\begin{tikzpicture}
\draw[world](3.5,2.4)--(3.5,1.95);\draw[world](3.5,0.35)arc(-90:90:.2);\draw[world](3.5,0.75)arc(-90:90:.2);\draw[world](3.5,1.15)arc(-90:90:.2);\draw[world](3.5,1.5500000000000003)arc(-90:90:.2);\draw[world](3.5,.35)--(3.5,0);
\draw[world](0,0)--(3.5,0) (0,2.4)--(1,2.4) (1.8,2.4)--(3.5,2.4) (0,0)--(0,.9) (0,1.3)--(0,2.4);\draw[world](1,2.2)rectangle(1.8,2.6);\node[above]at(1.4,2.65){$R$};\draw[world](-.45,.9)--(.45,.9) (-.45,1.3)--(.45,1.3);\node[left]at(-.55,1.1){$C$};\node[right]at(.5,1.5){$+q$};\node[right]at(.5,.6){$-q$};\node[right]at(3.9,1.1){$L$};\draw[flow,red](2,2.4)--(3,2.4)node[midway,above]{$i$};

\end{tikzpicture}}
图\ref{em5:rlc-circuit}与理想 LC 采用相同正向，先列 $L\dot i+Ri=q/C$，再代 $i=-\dot q$，便得到电荷方程。这样可以核对阻尼项的正号；若错误地得到 $-R\dot q$，就会预言无外源时能量自发增长。


含电阻的串联 RLC 自由电路，电荷满足 $L\ddot q+R\dot q+q/C=0$。设 $q=\e^{st}$，特征方程 $Ls^2+Rs+1/C=0$。当 $R<2\sqrt{L/C}$，根为 $-R/(2L)\pm\mathrm j\omega_d$，$\omega_d=\sqrt{1/(LC)-R^2/(4L^2)}$，振幅指数衰减；等号时临界阻尼；更大时两根均为负实数，没有周期反向的振荡。把所有含电感和电容的电路都称为等幅振荡是不对的。

阻尼电路的总能量变化率为 $-Ri^2$，始终非正，但瞬时衰减率不是常量。弱阻尼时才可将一个周期内的能量变化近似平均，得到包络按 $\e^{-Rt/L}$ 衰减。求振幅、周期或热量时，应区分严格解与弱阻尼近似，并检查电压和电流的连续初值是否都已使用。


\subsection{从相平面与时间图理解阻尼}

\fig{em5:phase-plane}{归一化电荷、电流相平面：理想回路的闭合轨道与欠阻尼回路的收缩轨道}{
\begin{tikzpicture}[scale=.88]
\begin{scope}
\draw[axis](-2.1,0)--(2.2,0)node[right]{$q/Q_*$};
\draw[axis](0,-2)--(0,2.2)node[above]{$i/(\omega_0Q_*)$};
\draw[world,blue](0,0)circle(1.5);
\draw[flow,blue](1.5,0)arc(0:35:1.5);
\draw[flow,blue](-1.5,0)arc(180:215:1.5);
\fill(1.5,0)circle(2pt);\node at(1.7,-1.15){$t=0$};
\node at(0,-2.6){$R=0$：能量保持};
\end{scope}
\begin{scope}[xshift=6cm]
\draw[axis](-2.1,0)--(2.2,0)node[right]{$q/Q_*$};
\draw[axis](0,-2)--(0,2.2)node[above]{$i/(\omega_0Q_*)$};
\draw[world,red,domain=0:28,samples=240,variable=\t]
 plot({1.5*exp(-.12*\t)*(cos(sqrt(.9856)*\t r)+.12/sqrt(.9856)*sin(sqrt(.9856)*\t r))},
 {1.5*exp(-.12*\t)/sqrt(.9856)*sin(sqrt(.9856)*\t r)});
\draw[flow,red](1.49,.08)--(1.45,.36);
\fill(1.5,0)circle(2pt);\node at(1.7,-1.15){$t=0$};
\node at(0,-2.6){$0<R<2\sqrt{L/C}$：能量耗散};
\end{scope}
\end{tikzpicture}}

图\ref{em5:phase-plane}取初态$q=Q_*>0,i=0$，并使用前文的放电正向$i=-\dot q$。起点处$\dot i=q/(LC)>0$，状态点先向上；到达上半平面后$\dot q<0$，所以沿逆时针方向前进。左图的归一化使能量椭圆变成圆；在未归一化的$q,i$坐标中仍是椭圆。右图采用同样初态，电阻使每次绕转后能量减小，轨道向原点收缩。两图的坐标均为状态量，不代表元件在空间中的位置。


理想LC的能量等值线$q^2/(2C)+Li^2/2=W_0$可与图\ref{em5:phase-plane}的运动方向结合使用：给定$q$时，能量只确定$|i|$，而正负号取决于状态点位于上半轨道还是下半轨道。电荷端点处$i=0$，电流端点处$q=0$，对应两种储能完全转换的时刻。

含电阻时，能量等值线不再保持，状态轨迹向原点收缩。欠阻尼时轨迹绕原点旋转，电容反复换极；临界与过阻尼时一般不出现这种周期绕转。对临界阻尼 $R=2\sqrt{L/C}$，特征根重合为 $-\alpha$，令 $q=y(t)\e^{-\alpha t}$，则
\[
\dot q=(\dot y-\alpha y)\e^{-\alpha t},\qquad
\ddot q=(\ddot y-2\alpha\dot y+\alpha^2y)\e^{-\alpha t}.
\]
代入 $\ddot q+2\alpha\dot q+\alpha^2q=0$ 后只余 $\ddot y=0$。先积得 $\dot y=B$，再积得 $y=A+Bt$，故通解必须写为 $(A+Bt)\e^{-\alpha t}$，不能把欠阻尼公式中的 $\omega_d$ 直接设为零而保留分母。

初电荷 $Q_0$、初电流为零时，$q(0)=A=Q_0$，$\dot q(0)=B-\alpha A=0$，解得 $B=\alpha Q_0$，临界解为 $q=Q_0(1+\alpha t)\e^{-\alpha t}$。过阻尼的两个负实根为 $s_1,s_2$，写 $q=A\e^{s_1t}+B\e^{s_2t}$，由 $A+B=Q_0$、$s_1A+s_2B=0$ 解出 $A=Q_0s_2/(s_2-s_1)$、$B=-Q_0s_1/(s_2-s_1)$，得 $q=Q_0(s_2\e^{s_1t}-s_1\e^{s_2t})/(s_2-s_1)$。这说明初态相同、参数不同的电路可以完全没有正弦式振荡，不能只根据含有 L 和 C 就画正弦波。

弱阻尼近似中，电荷振幅以 $\e^{-Rt/(2L)}$ 衰减，能量包络以 $\e^{-Rt/L}$ 衰减，两者指数相差一倍。用某两次同极性电荷峰值的比可以估计阻尼，再求能量损失；若直接把振幅比当能量比，会把损失估计错。全过程最终发热可由初总能量求得，但任意中间时间的热量还要扣除当前电场能与磁场能。

\fig{em5:damping-time}{相同初电荷、零初电流时，三种阻尼的电荷时间曲线}{
\begin{tikzpicture}[x=.9cm,y=2cm]
\draw[axis](0,0)--(9.5,0)node[right]{$\omega_0t$};
\draw[axis](0,-.7)--(0,1.25)node[above]{$q/Q_0$};
\draw[world,blue,domain=0:9,samples=160]plot(\x,{exp(-.2*\x)*(cos(sqrt(.96)*\x r)+.2/sqrt(.96)*sin(sqrt(.96)*\x r))});
\draw[world,red,dashed,domain=0:9,samples=100]plot(\x,{(1+\x)*exp(-\x)});
\draw[world,black,dash dot,domain=0:9,samples=100]
 plot(\x,{((2+sqrt(3))*exp((-2+sqrt(3))*\x)-(2-sqrt(3))*exp((-2-sqrt(3))*\x))/(2*sqrt(3))});
\node[left]at(0,1){$1$};
\node[blue]at(1.3,-1){欠阻尼};
\node[red]at(4.4,-1){临界阻尼};
\node at(7.7,-1){过阻尼};
\end{tikzpicture}}

图\ref{em5:damping-time}统一使用无阻尼角频率$\omega_0=1/\sqrt{LC}$量度时间。三条曲线分别取阻尼比$R/(2\sqrt{L/C})$为$0.2,1,2$，起点高度和初始斜率相同。蓝线穿过横轴，表示电容换极；另外两线在本图初态下保持正电荷并渐近衰减。不同初电流可能改变过零情况，因而不能把这一特定初态的图推广为任意初态都不换极。临界与过阻尼的判定依据仍是特征根及轨道性质。



\section{电磁机械振荡与两个连续状态}

零电阻多杆回路中，自感阻止电流突变，使机械相对位移与磁能之间可以周期交换。设等宽两杆的约化质量为 $\mu$、相对位移变化为 $\xi$，以初态间距为位移参考，并取初始电流为零。令 $k=Bl$，向纸里为正磁通，总磁链守恒给出 $Li+k\xi=0$。相对受力为 $\mu\ddot\xi=ki$，于是 $\ddot\xi+k^2\xi/(\mu L)=0$。频率取决于磁场耦合与约化质量，不能直接使用电容回路的 $1/\sqrt{LC}$，除非先建立真实等效关系。


\fig{em5:two-rods-inductance}{间距增加后实际电流逆时针，两棒受力使间距减小；质心运动另行保留}{\begin{tikzpicture}

\draw[world](0,0)--(4,0) (0,2.2)--(1.2,2.2) (2.8,2.2)--(4,2.2);\foreach\n in {0,...,3}{\draw[world]({1.2+.4*\n},2.2)arc(180:0:.2);}
\draw[world,very thick](0,0)--(0,2.2) (4,0)--(4,2.2);\node at(2,2.9){$L,\ R=0$};\node at(2,1.1){$\otimes B$};
\draw[flow,red](0,1.7)--(0,1)node[left]{$-i>0$};\draw[flow,red](4,.5)--(4,1.2)node[right]{$-i>0$};
\draw[force,blue](0,.55)--(1.1,.55)node[midway,below]{$-ki$};\draw[force,blue](4,1.75)--(2.9,1.75)node[midway,above]{$ki$};
\node[above]at(0,2.4){$m_1$};\node[above]at(4,2.4){$m_2$};\draw[dim](0,-.8)--(4,-.8)node[midway,below]{$x_2-x_1=x_0+\xi,\quad\xi>0$};

\end{tikzpicture}}
图\ref{em5:two-rods-inductance}用顺时针为参考电流正向：右棒中正电流向下，正磁力便向右。间距增大时 $i=-k\xi/L<0$，所以右棒受力向左、左棒受力向右，构成相对运动的回复作用。两棒相加仍有总动量守恒，不能将质心动能一并变为磁能。

相对运动的角频率与能量为
\[\Omega=\frac{k}{\sqrt{\mu L}},\qquad
\frac12\mu\dot\xi^2+\frac12Li^2
=\frac12\mu\dot\xi^2+\frac{k^2\xi^2}{2L}=\text{常量}.\]
固定 $L$ 是几何近似，要求棒间距改变不显著影响回路自感；若两棒接近到碰撞或交叉，原模型亦须终止。这里磁能对应相对位移的二次势能，与纯电阻回路的速度阻尼不同。


\section{受迫响应与自然响应怎样共同满足初态}

线性电路的完整响应可写成一个特解与齐次解之和。特解反映持续外源，自然响应负责使总电流、电压满足切换瞬间的状态条件。正弦电源接通后，最终正弦稳态只是特解；若接通相位不合适，初期还会叠加指数分量。电源相位相同不等于电流初相位相同，电感和电阻决定的相位滞后必须先算出。

二阶电路需要两个独立初始条件。以电容电荷为变量时，电感初始电流通过电荷导数进入；以导棒速度为变量时，电流初值通过加速度进入。求解以后，至少代回原一阶方程检查一次，以免消元时丢失符号。对于恒力驱动且含电感的导棒，速度可能围绕终值作衰减振荡，磁场力也可能短时间向运动供能，这只是磁能释放，并不违背总能量守恒。

峰值时刻应由所求电流或电压的导数为零确定；有外源或阻尼时，不能直接套用无阻尼四分之一周期。

\fig{em5:mechanical-state}{含自感导棒具有速度和电流两个状态；动生电动势与安培力形成相反的功率交换}{\begin{tikzpicture}

\draw[world](0,0)--(4,0) (0,2.2)--(1.2,2.2) (2.8,2.2)--(4,2.2) (0,0)--(0,.8) (0,1.5)--(0,2.2);
\foreach\n in {0,...,3}{\draw[world]({1.2+.4*\n},2.2)arc(180:0:.2);}\node at(2,2.8){$L$};\draw[world](-.2,.8)rectangle(.2,1.5);\node[left]at(-.3,1.15){$R$};
\draw[world,very thick](3.5,0)--(3.5,2.2);\node at(1.7,1){$\otimes B$};\draw[flow,red](3.5,.3)--(3.5,1)node[right]{$i$};\draw[flow](3.5,3.2)--(4.6,3.2)node[right]{$v$};
\draw[force](3.5,1.7)--(4.8,1.7)node[right]{$F$};\draw[force,blue](3.5,1.7)--(2.2,1.7)node[midway,above]{$ki$};\node at(1.9,-.7){$m\dot v=F-ki,\quad L\dot i+Ri=kv$};

\end{tikzpicture}}
图\ref{em5:mechanical-state}的两个一阶方程共同决定状态，不能只用 $Ri=kv$ 立即消去电流。消元得
\[mL\ddot v+mR\dot v+k^2v=RF,\qquad
v(0)=v_0,\quad \dot v(0)=\frac{F-ki_0}{m}.\]
初始电流作为初始加速度进入二阶方程。若 $R>0$，终态为 $i_\infty=F/k,v_\infty=RF/k^2$；欠阻尼条件为 $R<2k\sqrt{L/m}$。方程的两个状态还决定瞬时能量平衡
\[Fv=\frac{\dd}{\dd t}\left(\frac12mv^2+\frac12Li^2\right)+Ri^2.\]
电流有滞后时，不能仅凭速度正在下降就断言机械能已全部耗散；磁场也在储能或返还能量。


\fig{fv:f5:responses}{正弦电压从过零点接通时，正弦特解与自然项相加才满足零初流}{\begin{tikzpicture}

\draw[axis](-.2,0)--(13.1,0)node[right]{$\omega t$};\draw[axis](0,-1.5)--(0,2.35)node[above]{$i/I_m$};
\draw[blue,dashed,domain=0:12.57,samples=160]plot(\x,{sin(\x r-atan(5))});\draw[red,dashed,domain=0:12.57,samples=100]plot(\x,{sin(atan(5))*exp(-.2*\x)});
\draw[world,domain=0:12.57,samples=180]plot(\x,{sin(\x r-atan(5))+sin(atan(5))*exp(-.2*\x)});
\node[left]at(0,1){$1$};\node[left]at(0,-1){$-1$};\node[blue]at(9,-1.5){正弦特解};\node[red]at(8,1.85){自然项};\node at(4,2.35){完整响应};

\end{tikzpicture}}

固定参数的串联 RL 电路受正弦电压驱动时，设 $i_{\rm ss}(t)$ 是已求得的正弦特解，任意初流 $I_0$ 对应的完整解为
\[i(t)=i_{\rm ss}(t)+[I_0-i_{\rm ss}(0)]\e^{-t/\tau},\qquad\tau=L/R.\]
齐次项的系数由初始差额确定，不能在接通瞬间直接舍弃。图\ref{fv:f5:responses}取 $e=E_m\sin\omega t$、$R/(\omega L)=0.2$，稳态相位滞后 $\phi=\arctan5$，因此 $i_{\rm ss}(0)=-I_m\sin\phi$。红色自然项在零时刻恰抵消蓝色特解，使总电流为零；两者斜率也在起点抵消，与原方程给出的 $\dot i(0)=e(0)/L=0$ 一致。

若接通时已满足 $I_0=i_{\rm ss}(0)$，则没有指数暂态；这取决于电流的特解相位，不能简单用电压是否过零判断。电阻越小，自然项衰减越慢；理想纯电感的常量积分偏置不会自行消失。



\section{多储能元件中的独立振荡与剩余能量}

两个储能回路耦合以后，一般不能各用自己的 $1/\sqrt{LC}$ 描述。对称结构可把电荷或电流分成相加、相减的两个组合，使耦合方程分离。每种组合对应一种固有频率，实际运动是它们的叠加；弱耦合时两频率接近，叠加产生缓慢的能量交换。电荷振幅的包络与总能量的快振荡不是同一个函数，判断能量转移要同时计算电场、磁场和互感项。

\fig{em5:coupled-loops}{两个独立闭合回路只作磁耦合；两电流均流入同名端并从各自的带电极板流出}{\begin{tikzpicture}
\draw[world](3,2.4)--(3,1.95);\draw[world](3,0.35)arc(-90:90:.2);\draw[world](3,0.75)arc(-90:90:.2);\draw[world](3,1.15)arc(-90:90:.2);\draw[world](3,1.5500000000000003)arc(-90:90:.2);\draw[world](3,.35)--(3,0);\draw[world](6.3,2.4)--(6.3,1.95);\draw[world](6.3,0.35)arc(-90:90:.2);\draw[world](6.3,0.75)arc(-90:90:.2);\draw[world](6.3,1.15)arc(-90:90:.2);\draw[world](6.3,1.5500000000000003)arc(-90:90:.2);\draw[world](6.3,.35)--(6.3,0);
\foreach\x/\j in {0/1,9.3/2}{\draw[world](\x,0)--(\x,.9) (\x,1.3)--(\x,2.4);\draw[world]({\x-.4},.9)--({\x+.4},.9) ({\x-.4},1.3)--({\x+.4},1.3);\node[above]at(\x,1.5){$+q_{\j}$};\node[below]at(\x,.7){$-q_{\j}$};\node at(\x,-.6){$C$};}
\draw[world](0,0)--(3,0) (0,2.4)--(3,2.4) (6.3,0)--(9.3,0) (6.3,2.4)--(9.3,2.4);
\fill(2.75,2.15)circle(2pt);\fill(6.05,2.15)circle(2pt);\node[left]at(3,1.1){$L$};\node[right]at(6.7,1.1){$L$};\draw[<->](3.5,1.1)--(5.7,1.1)node[midway,above]{$M>0$};
\draw[flow,red](1,2.4)--(2.2,2.4)node[midway,above]{$i_1$};\draw[flow,red](8.5,2.4)--(7.3,2.4)node[midway,above]{$i_2$};

\end{tikzpicture}}
图\ref{em5:coupled-loops}采用 $i_j=-\dot q_j$。由互感端压分别列回路方程，得到
\[L\ddot q_1+M\ddot q_2+q_1/C=0,\qquad
M\ddot q_1+L\ddot q_2+q_2/C=0.\]
令 $q_\pm=q_1\pm q_2$，两式相加、相减便使方程分离：$(L\pm M)\ddot q_\pm+q_\pm/C=0$，频率为 $\omega_\pm=1/\sqrt{(L\pm M)C}$。每个独立组合均须使用自己的初值，最后再还原两回路的电荷与电流。磁场能中的交叉项 $Mi_1i_2$ 同样保留，不能只将两个单回路能量相加。


含电阻也不一定意味着所有初始磁能最终都变成热。两只理想电感并联时，若存在不经过电阻的环流组合，它可以持续存在；电阻只衰减与端口电压相联系的组合。末态由电流和为零以及一个磁链差不变量共同决定。先找不变量，再计算耗散部分，往往比逐支路强行积分更清晰。

\fig{em5:persistent-mode}{两个未耦合理想电感并联电阻释放能量；公共端口可以无流，电感间仍有环流}{\begin{tikzpicture}
\draw[world](0,2.4)--(0,1.95);\draw[world](0,0.35)arc(-90:90:.2);\draw[world](0,0.75)arc(-90:90:.2);\draw[world](0,1.15)arc(-90:90:.2);\draw[world](0,1.5500000000000003)arc(-90:90:.2);\draw[world](0,.35)--(0,0);\draw[world](3,2.4)--(3,1.95);\draw[world](3,0.35)arc(-90:90:.2);\draw[world](3,0.75)arc(-90:90:.2);\draw[world](3,1.15)arc(-90:90:.2);\draw[world](3,1.5500000000000003)arc(-90:90:.2);\draw[world](3,.35)--(3,0);
\draw[world](0,2.4)--(0,2.8)--(6,2.8)--(6,1.5) (0,0)--(0,-.4)--(6,-.4)--(6,.8) (3,2.4)--(3,2.8) (3,0)--(3,-.4);
\draw[world](5.8,.8)rectangle(6.2,1.5);\node[right]at(6.3,1.15){$R$};\node[left]at(0,1.1){$L_1$};\node[left]at(3,1.1){$L_2$};
\draw[flow,red](-.6,2.5)--(-.6,1.8)node[left]{$i_1$};\draw[flow,red](3.8,2.5)--(3.8,1.8)node[right]{$i_2$};\draw[flow,blue](6.7,2.5)--(6.7,1.8)node[right]{$i_R$};
\fill(4.8,2.8)circle(2pt)node[above]{$A$};\fill(4.8,-.4)circle(2pt)node[below]{$B$};\node at(3,-1.1){$i_R=-(i_1+i_2),\quad u=V_A-V_B$};

\end{tikzpicture}}
图\ref{em5:persistent-mode}中三条支路均取向下为正。节点关系给出 $i_R=-(i_1+i_2)$，共同端压给出 $L_1\dot i_1=L_2\dot i_2=u=-R(i_1+i_2)$。令 $S=i_1+i_2$，相加后
\[\dot S=-R\left(\frac1{L_1}+\frac1{L_2}\right)S,\qquad
L_1i_1-L_2i_2=D_0=\text{常量}.\]
令 $\tau=L_1L_2/[R(L_1+L_2)]$，总电流方程的分离变量式为
\[
\int_{S(0)}^{S(t)}\frac{\dd z}{z}=-\int_0^t\frac{\dd s}{\tau},
\qquad S(t)=S(0)\e^{-t/\tau}.
\]
若 $S(0)=0$，直接取 $S(t)=0$；磁链差则由
$\dd(L_1i_1-L_2i_2)/\dd t=u-u=0$ 得到，因此不衰减。末态 $S=0$，于是 $i_{1\infty}=D_0/(L_1+L_2)$、$i_{2\infty}=-i_{1\infty}$，可保留非零环流。完整磁能可分解为
\[W_B=\frac12\frac{L_1L_2}{L_1+L_2}S^2+
\frac{D_0^2}{2(L_1+L_2)}.\]
前一项最终耗散，后一项作为磁能保留。若要使全部电流消失，必须另有能耗散这一环流组合的实际电阻，不能凭网络中存在某个电阻就判定所有磁能必然转为热。


这些情况说明，判断最终状态需要同时考虑电路拓扑、守恒约束和耗散通道。只写“足够久以后电流为零”并不是普遍成立的物理结论。下面的例题分别处理有限脉冲、限压续流、接通相位、快切换、耦合振荡及机电耦合，保留每种模型真正不同的状态连接关系。

